\documentclass[11pt]{article}
\usepackage[margin=1in]{geometry}
\usepackage{amsmath,amssymb,mathtools}
\usepackage{booktabs}
\usepackage{enumitem}
\usepackage{siunitx}
\sisetup{group-separator={,},group-minimum-digits=3}

\title{ECON 101: Macroeconomics \\ Quiz 2 -- Solutions}
\author{Instructor: Muhebullah Karimzada}
\date{October 01,  2025}

\begin{document}
\maketitle



\section*{Part A — Multiple Choice (1 mark each)}

\paragraph{1.} \textbf{(c) Demand for chicken increases.}\\
\emph{Reasoning:} Beef and chicken are substitutes. A higher price of beef makes chicken relatively cheaper, which increases the \emph{demand} for chicken (the entire demand curve for chicken shifts right). It is not a movement along the chicken demand curve.

\paragraph{2.} \textbf{(b) A Canadian firm sells lumber to a U.S. company.}\\
\emph{Reasoning:} GDP records the market value of final goods and services produced \emph{within} Canada. Lumber produced in Canada and sold abroad is an \emph{export} and is included in Canada’s GDP. Used goods sales are excluded because they were counted when first produced. Spending by Canadians abroad is counted in the other country’s GDP. Home production is excluded.

\paragraph{3.} \textbf{(c) Values output at constant prices.}\\
\emph{Reasoning:} Real GDP measures quantities using a fixed base year’s prices. This removes the effect of price changes and isolates quantity changes. Nominal GDP uses current prices.

\paragraph{4.} \textbf{(a) 10\%.}\\
\emph{Definition:} Inflation (via GDP deflator) from period 1 to 2 is
\[
\pi=\frac{P_2-P_1}{P_1}\times 100.
\]
\emph{Calculation:}
\[
\pi=\frac{132-120}{120}\times 100=\frac{12}{120}\times 100=0.1\times 100=10\%.
\]

\paragraph{5.} \textbf{(c) An increase in unpaid household work.}\\
\emph{Reasoning:} Unpaid household work is not transacted in markets, so it is excluded from GDP. New consumption, higher exports, and higher investment are all counted and would increase GDP.

\bigskip
\hrule
\bigskip
\newpage

\section*{Part B — Numerical Problems (15 marks total)}

\subsection*{1. Market Equilibrium and a Price Floor}

\paragraph{Given:}
\[
Q_D=200-5P,\qquad Q_S=20+3P,
\]
where \(Q\) is quantity and \(P\) is price in dollars.

\subsubsection*{(a) Equilibrium price and quantity}
At equilibrium, quantity demanded equals quantity supplied:
\[
Q_D=Q_S.
\]
\emph{Set up and solve for \(P\):}
\[
200-5P=20+3P
\quad\Rightarrow\quad
200-20=3P+5P
\quad\Rightarrow\quad
180=8P
\quad\Rightarrow\quad
P^\ast=\frac{180}{8}=22.5.
\]
\emph{Find \(Q^\ast\) by substituting \(P^\ast\) into either curve:}
\[
Q_S=20+3(22.5)=20+67.5=87.5.
\]
Check with demand:
\[
Q_D=200-5(22.5)=200-112.5=87.5.
\]
\emph{\textbf{Answer:}} \(P^\ast=\$22.50\), \(Q^\ast=87.5\) units.

\subsubsection*{(b) Price floor at \$25: surplus or shortage}
A price floor sets a legal minimum price. It is \emph{binding} if it is above the equilibrium price. Here, \(25>22.5\), so it is binding.

\emph{Compute quantities at \(P=\$25\):}
\[
Q_D=200-5(25)=200-125=75,\qquad
Q_S=20+3(25)=20+75=95.
\]
\emph{Compare:} \(Q_S>Q_D\). The difference is a \textbf{surplus}:
\[
\text{Surplus}=Q_S-Q_D=95-75=20\ \text{units}.
\]
\emph{\textbf{Answer:}} Binding floor creates a surplus of \(20\) units.

\bigskip
\hrule
\bigskip
\newpage
\subsection*{Q7. Expenditure Approach to GDP}

\paragraph{Given (billions):}
\[
C=1{,}200,\quad I=500,\quad G=600,\quad X=450,\quad M=550.
\]
Expenditure GDP is
\[
Y=C+I+G+(X-M).
\]

\subsubsection*{(a) Nominal GDP}
\emph{Compute net exports:}\ \(\text{NX}=X-M=450-550=-100.\)\\
\emph{Compute GDP:}
\[
Y=1{,}200+500+600+(-100)=2{,}200\ \text{billion}.
\]
\emph{\textbf{Answer:}} \(\$2{,}200\) billion.

\subsubsection*{(b) GDP per capita}
 GDP per capita is GDP divided by population. Population \(=37\) million.
\[
\text{GDP per capita}=
\frac{2{,}200\ \text{billion}}{37\ \text{million}}
=
\frac{2{,}200\times 10^9}{37\times 10^6}
=
\frac{2{,}200}{37}\times 10^3
\approx 59{,}459.46.
\]
\emph{\textbf{Answer:}} \(\$59{,}459.46\) per person (approximately).

\subsubsection*{(c) Output gap as a percentage}
Output gap percentage is
\[
\text{Gap}=\frac{Y-Y^\ast}{Y^\ast}\times 100,
\]
where \(Y^\ast\) is potential GDP. Given \(Y^\ast=\$3{,}000\) billion.

\emph{Compute:}
\[
\frac{2{,}200-3{,}000}{3{,}000}\times 100
=
\frac{-800}{3{,}000}\times 100
=
-26.666\ldots\%
\approx -26.7\%.
\]
\emph{\textbf{Answer:}} \(-26.7\%\) recessionary gap.

\bigskip
\hrule
\bigskip

\subsection*{Q8. Real vs. Nominal GDP and the GDP Deflator}

\paragraph{Data:}
\begin{center}
\begin{tabular}{c c c c c}
\toprule
Year & Quantity Milk & Price Milk & Quantity Computers & Price Computers \\
\midrule
2023 & 1{,}000 & \$2 & 50 & \$500 \\
2024 & 1{,}200 & \$3 & 55 & \$550 \\
\bottomrule
\end{tabular}
\end{center}

\subsubsection*{(a) Nominal GDP for 2023 and 2024}
Nominal GDP uses current-year prices and quantities:
\[
\text{Nominal }Y_t=\sum_i p_{it}q_{it}.
\]
\emph{2023:}
\[
1{,}000\cdot 2\;+\;50\cdot 500
=2{,}000+25{,}000
=27{,}000.
\]
\emph{2024:}
\[
1{,}200\cdot 3\;+\;55\cdot 550
=3{,}600+30{,}250
=33{,}850.
\]
\emph{\textbf{Answers:}} 2023 = \(\$27{,}000\), 2024 = \(\$33{,}850\).

\subsubsection*{(b) Real GDP for 2024 using 2023 as base year}
Real GDP values quantities at base-year prices:
\[
\text{Real }Y_{2024}^{\text{(base 2023)}}=
p_{\text{milk},2023}\cdot q_{\text{milk},2024}
+
p_{\text{comp},2023}\cdot q_{\text{comp},2024}.
\]
\emph{Compute:}
\[
=2\cdot 1{,}200 + 500\cdot 55
=2{,}400+27{,}500
=29{,}900.
\]
\emph{\textbf{Answer:}} \(\$29{,}900\).

\subsubsection*{(c) GDP deflator for 2024}
GDP deflator is a price index:
\[
\text{Deflator}_{2024}=\frac{\text{Nominal }Y_{2024}}{\text{Real }Y_{2024}^{\text{(base 2023)}}}\times 100.
\]
\emph{Compute:}
\[
\frac{33{,}850}{29{,}900}\times 100
\approx 1.1321\times 100
\approx 113.21.
\]
\emph{\textbf{Answer:}} \(\approx 113.2\).

\subsubsection*{(d) Inflation rate from 2023 to 2024}
\emph{Note:} With base year 2023, \(\text{Deflator}_{2023}=100\).\\

\[
\pi=\frac{\text{Deflator}_{2024}-\text{Deflator}_{2023}}{\text{Deflator}_{2023}}\times 100.
\]
\emph{Compute:}
\[
\pi=\frac{113.21-100}{100}\times 100
=13.21\%.
\]
\emph{\textbf{Answer:}} \(\approx 13.2\%\).


\end{document}
